\documentclass[book.tex]{subfiles}
\begin{document}

\chapter{Polymers}

\label{chap:polymers}
\noindent\rule{11cm}{0.4pt} \\
\textbf{Prerequisites:}  \autoref{chap:phase_equilibrium}  \\
\textbf{Difficulty Level:} ** \\
\noindent\rule{11cm}{0.4pt}
\vspace{5mm}

The term polymer is used to describe molecules which consist of a large number of repeating units connected by covalent chemical bonds. Some polymers are linear, like beads on a necklace; others are branched. Statistical mechanics has a prominent role in predicting the properties of polymers. For example, stretching the polymer chains lowers their conformational entropy, and so under no external force conditions, the chain naturally tend to retract to regain entropy. The entropy is also lowered when polymers become compact, as when proteins fold or when DNA becomes encapsulated within virus heads. Additionally the changes on the end-to end length of polymer molecules due stretching, can be used to interpret the viscosities of polymer solutions, the scattering of light, and some of the dynamic properties of polymer solutions. 

In order to understand the thermodynamic behavior due to polymer stretching, a proper characterization of the random changes in length needs to be considered. For this, we will consider the \textit{random flight} model. This model has the advantage that in the case of long chains, it can be approximated as a Gaussian distribution, and so simplifying the analysis of the statistical mechanics.


\section{Polymers and proteins}

Polymer molecules are chains of covalently bonded monomer units (Fig.~\ref{fig1}).

\begin{marginfigure}
\includegraphics[width=\linewidth]{figures/Physics/statistical_mechanics/polymers/fig1.png}
\includegraphics[width=\linewidth]{figures/Physics/statistical_mechanics/polymers/fig2.png}
\caption{Some examples of Polymer. The figure above corresponds to the chemical structure of polyethylene. The figure below is the chemical structure of DNA}
\label{fig1}
\end{marginfigure}

Proteins are polymers that are composed of specific sequence of
amino acids that are linked by peptide bonds. Nature has exquisite control over primary structure of proteins (\textit{i.e.} amino acid sequence), facilitating the manufacture of proteins that fold into precise structures with specific biological functions. There are 20 natural amino acids that are used to construct biological proteins. The diverse chemical functionality of these amino acids provides the flexibility to create enzymes and structural proteins that engage in a wide range of biological processes (Fig.~\ref{fig8}).

%\begin{marginfigure}
%\includegraphics[width=\linewidth]{figures/Physics/statistical_mechanics/polymers/fig8.png}
%\caption{The twenty amino acids divided into six categories according to the chemical nature of their side chains.}
%\label{fig8}
%\end{marginfigure}

\section{Physical behavior at different length scales}

The chemical structure of a polymer dictates the atomic level physical behavior. Thermal fluctuations at small length scales result in local deformations that are randomly distributed about a mean structure with a finite variance. At larger length scales, the correlation between the fluctuating chain segments vanishes, and the behavior is that of many independently fluctuating chain segments. 

The Central Limit Theorem states that any sum of many independent identically distributed random variables will tend to a normal (or Gaussian) distribution, provided the sum of variables has a finite variance.

\begin{marginfigure}
\includegraphics[width=\linewidth]{figures/Physics/statistical_mechanics/polymers/fig10.png}
\caption{A polyethylene molecule at different length scales. The
atomic length scale exhibits a very specific physical behavior. However, the specifics of the underlying physics is washed out as we proceed to larger and larger length scales.}
\label{fig10}
\end{marginfigure}

Various models for the small-length-scale physical behavior of a polymer exist (see Fig.4). Though they differ in their small length scale behavior, they all tend to a Gaussian distribution at large length scales. We focus on a discrete random walk (or a freely jointed chain) to demonstrate the crossover to a Gaussian distribution. 

\begin{figure}
\includegraphics[width=\linewidth]{figures/Physics/statistical_mechanics/polymers/fig11.png}
\stepcounter{figure}
\smallskip\noindent\small Figure \thefigure: Several ideal chain models.
\label{fig11}
\end{figure}

The Ideal random-walk (or freely jointed) polymer chain is the simplest idealization of a flexible polymer. The polymer does not interact with itself; we can think of it as a phantom chain. This framework is mathematically useful since random walks appear in a range of seemingly different physical processes (diffusion, heat transfer, quantum mechanics, etc).


\section{Entropic elasticity}

Consider a walk on a one-dimensional lattice. After $N$ steps, how far do you travel on average?

\begin{align}\label{Eq1}
X = s_1+s_2+...+s_N=\sum_{i=1}^Ns_i
\end{align}

\noindent where $s_i = \pm 1$ is the jump vector. The mean-square end displacement is:
\begin{align}\label{Eq2}
\langle X^2\rangle = \sum_{i=1}^N\sum_{j=1}^N\langle s_is_j\rangle
\end{align}

Since jumps are uncorrelated, we have:
\begin{align}\label{Eq3}
\langle s_i s_j \rangle = \delta_{ij}
\end{align}

\noindent hwere $\delta_{ij}$ if $i=j$ and $\delta_{ij}=0$ if $i\ne j$. Therefore
\begin{align}\label{Eq4}
\langle X^2\rangle = \sum_{i=1}^N\sum_{j=1}^N \delta_{ij} = \sum_{i=1}^N 1 =N
\end{align}
This gives the size of a 1-dimensional polymer with discrete subunits of unit length. 

This is extended to 3 dimensions. Consider a chain of $N$ bond
vectors $\vec{b}_i$ ($i =1,2,...,N$) with length $b (\vec{b}_i \cdot\vec{b}_i =b^2)$, with end-to-end vector $\vec{R} = \sum_{i=1}^N\vec{b}_i$. The mean-square end-to-end distance is:
\begin{align}\label{Eq5}
\langle \vec{R}^2 \rangle = \sum_{i=1}^N\sum_{j=1}^N\langle\vec{b}_i\vec{b}_j\rangle = \sum_{i=1}^N\sum_{j=1}^N b^2\delta_{ij}=b^2N
\end{align}

\noindent using similar arguments as the 1-dimensional case.

Decomposing the end-to-end vector into its components $\vec{R} = X\hat{x}+ Y\hat{y}+Z\hat{z}$, we have:
\begin{align}\label{Eq6}
\langle X^2 \rangle = \langle Y^2 \rangle = \langle Z^2 \rangle = \frac{b^2N}{3}
\end{align}

\noindent since no direction is preferred.

The central limit theorem states that the limiting behavior of any statistical distribution tends to a Gaussian distribution in the limit of large sample size. In our case, the limit of large $N$ leads to a chain probability distribution that is Gaussian. With this, we construct the polymer chain distribution function.

Generally, a Gaussian distribution for a variable $X$ with zero
mean ($\langle X \rangle = 0$) and variance $\langle X^2 \rangle$
is written as:
\begin{align}\label{Eq7}
p_x(X,N)=\frac{1}{\sqrt{2\pi\langle X^2\rangle}}\exp\left( -\frac{X^2}{2\langle X^2 \rangle} \right)
\end{align}

This is the limiting behavior of a random walk in one dimension
as the number of steps is very large ($N \gg 1$).

Assuming the probability distribution is decoupled in the 3 directions, we write the 3-dimensional probability distribution as
\begin{align}\label{Eq8}
p(\vec{R},N) &= p_x(X,N)p_y(Y,N)p_y(Z,N) \notag\\
&=\left(\frac{2\pi Nb^2}{3}\right)^{-3/2}\exp\left( -\frac{3\vec{R^2}}{2Nb^2} \right)
\end{align}

\noindent giving the probability that a chain of length $N$ that begins at the origin will end at position $\vec{R}$.

The chain probability gives the free energy for constraining the
chain ends at a fixed position; specifically, we have:
\begin{align}\label{Eq9}
F(\vec{R})-F(\vec{0}) = -k_BT\log\left[ \frac{p(\vec{R},N)}{p(\vec{0},N)} \right] = \frac{3k_BT}{2Nb^2}\vec{R}^2
\end{align}

Thus, the chain behaves as a Hookean spring, \textit{i.e} a Gaussian
chain exhibits a linear response to tension. Pulling a Gaussian chain along the $z$-axis a distance $Z$ requires a tension $\tau$ given by:
\begin{align}\label{Eq10}
\tau = \frac{\partial F(Z)}{\partial Z} = \frac{3k_BT}{Nb^2}Z
\end{align}

\begin{marginfigure}
\includegraphics[width=\linewidth]{figures/Physics/statistical_mechanics/polymers/fig12.png}
\caption{The response of DNA to tension \citep{Smith1992} compared to three competing theories: the Gaussian chain (dotted green curve), the freely jointed chain (dashed blue curve), and the wormlike chain (red dashed-dotted curve are asymptotic results and solid red curve are exact results from Ref.\citep{Spakowitz2004}). In all cases, the Kuhn length is set to $b = 2l_p = 106nm$.}
\label{fig12}
\end{marginfigure}

This linear response governs the small-scale deformation of a polymer chain. The response is purely entropic, indicated by the linear scaling in temperature. This leads to the curious behavior for a stretch elastic band retracts when heated. The Gaussian chain model does not capture the behavior once the chain is nearly extended. Full extension reflects the microscopic details of the model; for example, the nearly inextensible covalent bonds of the polymer chain. However, many polymer properties are captured by the Gaussian chain model. This is highly advantageous because the properties tend to be governed by just a few parameters (Fig.~\ref{fig12}).

\section{Freely Jointed Model}

Consider one link of a polymer chain shown in the Fig.~\ref{fjm} which has a fixed length of $b$. The bond angle is given as $\theta$ and the torsion angle is give as $\phi$, both being free to to rotate. A force of magnitude $f$ is applied at its ends. Here, we will develop a model for its streching.

\begin{marginfigure}
\includegraphics[width=\linewidth]{figures/Physics/statistical_mechanics/polymers/fjm.png}
\caption{A polymer with a few links represented in the spherical coordinates. The force $\vec{f}$ denoted by arrow acts in the $\hat{z}$ direction as shown. The angle of a general link, represented by $\vec{r} = b\hat{t}$, on the \textit{z} axis is $\theta$.}
\label{fjm}
\end{marginfigure}

Since each polymer link has a fixed length $b$, the system has only 2 degrees of freedom. When this single link is pulled with a force, $\vec{f} = f \hat{z} $, 
the energy associated with it is given by $\varepsilon = -\vec{r}.\vec{f}$. The partition function $q(f)$ would be:

\begin{align}
q &= \sum \exp(-\beta E) \\
&= \int_0^{2\pi}\int_0^\pi\exp[-\beta(-\vec{r}\cdot \vec{f})]\sin\theta d\theta d\phi\\
&=\int_0^{2\pi}\int_0^\pi\exp(\beta zf)\sin\theta d\theta d\phi\\
&=\frac{2\pi}{b}\int_{-b}^{+b}\exp(\beta zf) dz\\
&=\frac{4\pi}{bf}k_BT\sinh\left(\frac{fb}{k_BT}\right)
\end{align}

Consider now a polymer made of $N$ such links. Since the links are distinguishable, we can write down the partition function $Q(f,N)$ as 
 $$Q = q^N$$
 
Using the above partition function, we can evaluate the Helmholtz free energy, $F(f,N)$ as: 

\begin{align}
F = -k_BT\log Q 
\end{align}

The Helmholtz free energy can be used to calculate the average displacement $\langle r \rangle$.

Since $F = -\langle r\rangle f \rightarrow \langle r \rangle = \frac{\partial F}{\partial f}$, we have:
\begin{align*}
\langle r\rangle  &= k_BT\frac{\partial \log Q}{\partial f} = \frac{k_BT}{q^N}Nq^{N-1}\frac{\partial q}{\partial f} \\
&= \frac{k_BTN}{q}\left[ \frac{4\pi}{f} \cosh\left( \frac{fb}{k_BT} \right) - \frac{4\pi}{bf^2} \sinh\left( \frac{fb}{k_BT} \right) \right] \\
&= \frac{N}{f}\left[ bf\coth\left( \frac{fb}{k_BT} \right) - k_BT \right]
\end{align*}


\begin{marginfigure}
\includegraphics[width=\linewidth]{figures/Physics/statistical_mechanics/polymers/stretch.jpg}
\caption{The average displacement $\langle r \rangle$ as a function of $f$  in the range 1 to 100 pN (pN = 10$^{-12}$ N), given $k_BT$ = 0.0259 eV, b = 0.5 nm and N = 20.}
\label{stretch}
\end{marginfigure}

\end{document}